% MAGVIT-v2 编码器的时间因果下采样（核长3、步长2、向过去补2帧），只画时间维；宽约100mm。
\begin{tikzpicture}[font=\figurefont\small,x=4.3mm,y=1cm,
  frame/.style={draw=FigureBlue,fill=FigureBlueFill,line width=0.6pt,
    minimum width=3.5mm,minimum height=3.5mm,inner sep=0pt},
  mid/.style={frame,draw=FigureYellow,fill=FigureYellowFill},
  top/.style={frame,draw=FigureRed,fill=FigureRedFill},
  pad/.style={frame,draw=BookMuted,fill=black!8},
  link/.style={draw=BookMuted,line width=0.5pt},
  note/.style={text=BookMuted,font=\figurefont\footnotesize,anchor=east}]
  % 第一行：17 帧输入，前补 2 帧
  \node[note] at (-3.7,0) {17 帧输入};
  \node[pad] at (-1,0) {};
  \node[pad] at (0,0) {};
  \foreach \i in {1,...,17}{\node[frame] (a\i) at (\i,0) {};}
  \node[font=\figurefont\footnotesize] at (1,0.45) {1};
  \node[font=\figurefont\footnotesize] at (17,0.45) {17};
  % 第二行：9 个位置，与各自窗口中最新的一帧对齐，前补 2 个位置
  \node[note] at (-3.7,-1.3) {9 个位置};
  \node[pad] at (-3,-1.3) {};
  \node[pad] at (-1,-1.3) {};
  \foreach \j in {1,...,9}{
    \pgfmathtruncatemacro{\x}{2*\j-1}
    \node[mid] (b\j) at (\x,-1.3) {};}
  % 第三行：5 个位置
  \node[note] at (-3.7,-2.6) {5 个位置};
  \foreach \k in {1,...,5}{
    \pgfmathtruncatemacro{\x}{4*\k-3}
    \node[top] (c\k) at (\x,-2.6) {};}
  % 示例窗口：第二行第 3 个位置读第 3、4、5 帧；第三行第 2 个位置读第二行第 1、2、3 个位置
  \foreach \i in {3,4,5}{\draw[link] (a\i.south) -- (b3.north);}
  \foreach \j in {1,2,3}{\draw[link] (b\j.south) -- (c2.north);}
  % 第一个位置只读补位与第 1 帧
  \draw[link,dashed] (-1,-0.18) -- (b1.north);
  \draw[link,dashed] (0,-0.18) -- (b1.north);
  \draw[link] (a1.south) -- (b1.north);
  \draw[link,dashed] (-3,-1.48) -- (c1.north);
  \draw[link,dashed] (-1,-1.48) -- (c1.north);
  \draw[link] (b1.south) -- (c1.north);
\end{tikzpicture}
